%! TeX root = ../thesis.tex
\chapter{Implementation}

Currently,
the limiting factor is the lookup of bit patterns inside the dictionary.
We therefore want to reduce the number of dictionary keys,
which we implement under the name ``hybrid representation'' (HR).

\section{Hybrid wave function}

For the geometry shown in~\cref{fig:wave-function-split}
we can define two sectors.
The first contains all sites on which $\ham_s$ acts,
the second contains the sites in the two chains of $\ham_v$.
Given a fixed number of electrons
we can distribute them over these sectors such that
for each sector we obtain a set of electronic configurations.
Every single determinant state can be expressed
by combining one configuration from each sector:
$\ket{s_i}$ from $\ham_s$ and $\ket{v_i}$ from $\ham_v$.
The first sector uses the existing bit representation (BR)
to denote a configuration.
In the second sector however,
we map the configurations onto linearly independent vectors
$\ket{v_i} = \sum_\alpha v_{i\alpha} \ket{\alpha}$
with each index $\alpha$ corresponding to a given state $\ket{\alpha}$
of the truncated chains
and the value $v_{i\alpha}$ corresponding to the probability amplitude.
This way every state of the full system can be expressed as
\begin{equation}
    \ket{\psi} = \sum_i \ket{s_i} \otimes \ket{v_i}
    = \sum_{i, \alpha} v_{i\alpha} \ket{s_i} \otimes \ket{\alpha}.
\end{equation}
Normalization is fulfilled by
\begin{equation}
    \sum_i \|\ket{v_i}\|^2=\sum_{i, \alpha} |v_{i\alpha}|^2 = 1.
\end{equation}
For a system depicted in~\cref{fig:wave-function-split} limited to singles,
and starting with 3 empty and 3 filled states in $\ham_v$,
and a factor 2 for spin,
the vector dimension would be $1+2\cdot(3+3) = 13$;
one entry stores the amplitude for no excitation
and 12 entries for an excitation in any of the possible spin-sites.

For this, we can use the struct \mintil{CIWavefunction}.
\begin{minted}[frame=lines]{julia}
struct CIWavefunction{T,nslater,nfilled,nempty,excitation}
    terms::T
end
\end{minted}
Determinants with their corresponding vectors are stored
inside \mintil{terms} as a dictionary.
To uniquely describe this new wave function,
we must save the number of sites inside the determinant $n_s$,
number of filled (empty) sites $n_f$ ($n_e$) inside the vector,
and maximum possible CI excitations $e$.
In this way, the structure is defined at compile time rather than
runtime if these values were field names instead.
We can define the interface as:
\begin{itemize}
    \item \mintil{terms()} returning the dictionary,
    \item \mintil{nslater()} returning $n_s$,
    \item \mintil{nfilled()} returning $n_f$,
    \item \mintil{nempty()} returning  $n_e$,
    \item \mintil{excitation()} returning $e$,
    \item \mintil{nvector()} returning $n_v = n_f + n_e$,
    \item \mintil{nsites()} returning $n_s + n_f + n_e$.
\end{itemize}

\subsection{Vector indexing}
\label{subsec:vector-indexing}

For vector indexing we must establish a bijective mapping between
indices and states.
We could group the spin-sites inside the SD either by spin
\begin{equation}
    \label{eq:sort-by-spin}
    \ket{
        \underbrace{
            \ldots, c_{2\up}, c_{1\up}, \ldots, v_{2\up}, v_{1\up}, s_\up
        }_{\up},
        \underbrace{
            \ldots, c_{2\dn}, c_{1\dn}, \ldots, v_{2\dn}, v_{1\dn}, s_\dn
        }_{\dn}
    },
\end{equation}
or by site
\begin{equation}
    \label{eq:sort-by-site}
    \ket{
        \ldots,
        \underbrace{c_{2\up}, c_{2\dn}}_{c_2},
        \underbrace{c_{1\up}, c_{1\dn}}_{c_1},
        \ldots,
        \underbrace{v_{2\up}, v_{2\dn}}_{v_2},
        \underbrace{v_{1\up}, v_{1\dn}}_{v_1},
        s
    }.
\end{equation}
To decide between these two possibilities,
we can calculate the effect of a hopping term $t_{ij} \cdag_i \cc_j$ on a SD
using the anticommutator definition given in
\cref{eq:anticommutator-definition}.
For $i > j$ we get
\begin{equation}
    \label{eq:hopping-1}
    \begin{split}
         & t_{ij} \cdag_i \cc_j \ket{\ldots, n_i=0, \ldots, n_j=1, \ldots} \\
         & = t_{ij} \prod_{k=j+1}^{i-1} (-1)^{n_k}
        \ket{\ldots, n_i=1, \ldots, n_j=0, \ldots},
    \end{split}
\end{equation}
and for $i < j$
\begin{equation}
    \label{eq:hopping-2}
    \begin{split}
         & t_{ij} \cdag_i \cc_j \ket{\ldots, n_j=1, \ldots, n_i=0, \ldots} \\
         & = t_{ij} \prod_{k=i+1}^{j-1} (-1)^{n_k}
        \ket{\ldots, n_j=0, \ldots, n_i=1, \ldots},
    \end{split}
\end{equation}
As we can see,
the sign of the result depends on the number of occupied
sites between $i$ and $j$.
Since our chains only contain nearest-neighbor hopping,
we can group by spin and skip the sign calculation.

Let the superscript denote the excitation and the subscript the index within
the vector.
We can store excitations sequentially,
e.g., if $\ket{v^0} = (v_1^0)$ contains no excitations,
$\ket{v^1} = (v_1^1, v_2^1, \ldots)$ all single excitations, etc.,
the full vector can then be assembled by listing the entries of each component
$\ket{v} = (v_1^0, v_1^1, v_2^1, \ldots)$.
The vector dimension is therefore
\begin{equation}
    \dim(\ket{v})
    = \sum_{i=0}^e \dim(\ket{v^i})
    = \sum_{i=0}^e \binom{2n_v}{e},
\end{equation}
where $n_v$ is the number of sites inside $\ham_v$
and the factor 2 accounts for spin.
Since $\dim(\ket{v^0}) = 1$, no index calculations are necessary for
no excitation.
For single excitation inside $\ket{v^1}$,
we must establish a mapping
\begin{equation}
    f: \{1, \ldots, n_v\} \times \{1, 2\} \rightarrow \{1, \ldots, 2n_v\},
\end{equation}
by assigning each spin-site
$(k, \sigma)\mapsto k + (\sigma - 1) n_v \eqcolon a$ an index $a$
and storing the amplitude of this excitation in $v_a^1$.
This function $f$ and its inverse are implemented under
\mintil{_rank1()} and \mintil{_unrank1()}.

For higher excitations we have multiple values $a_i$,
which we again want to map to an index.
For this,
we can use the combinatorial number system
(combinadics)~\cite{Knuth2005}.
Given a sequence of numbers $2n_v \geq a_k > \ldots > a_2 > a_1 \geq 0$,
it assigns a zero-based index $N$
by the sum of binomial coefficients
\begin{equation}
    N = \sum_{i=1}^{k} \binom{a_i}{i},
\end{equation}
with $\binom{n}{k} = 0$ for $k > n$.
This operation is called ``ranking'',
providing a one-to-one mapping between sequence of numbers and integers
without any gaps~\cite{Liang1995}.
The inverse operation (``unranking'') can be done iteratively.
In the first step we find the biggest integer $a_k$ which satisfies
\begin{equation}
    \label{eq:unranking}
    a_k = \max_{a_k \in \mathbb{N}}\left(\binom{a_k}{k} \leq N\right).
\end{equation}
The next number can be found by subtracting the result from the previous index.
\begin{equation}
    a_{k-1} = \max_{a_{k-1} \in \mathbb{N}}
    \left(\binom{a_{k-1}}{k-1} \leq N'\right), \qquad N' = N - \binom{a_k}{k}.
\end{equation}
For double excitation however,
we can use analytical expressions
\begin{align}
    a_2 & = \left\lfloor \frac{1 + \sqrt{1 + 8N}}{2} \right\rfloor, \\
    a_1 & = N - \binom{a_2}{2}.
\end{align}
Since Julia indexing is one-based, numbers must be shifted accordingly.
Mapping between spin-sites and index are available under the functions
\mintil{getvectorindex()} returning the vector index given excited spin-sites,
and \mintil{getspinsites()} returning excited spin-sites given a vector index.

\section{Hybrid operator}

To represent the Anderson Hamiltonian which acts on a hybrid wave function,
we can create a struct \mintil{CIOperator}.
As mentioned in~\cref{sec:AIM} it should contain three parts,
one operator acting on $\ket{s}$,
one acting on $\ket{v}$ using an expansion,
and a mixed operator representing hopping between
$\ket{s}$ and $\ket{v}$
\begin{equation}
    \label{eq:hamiltonian-split}
    \ham = \ham_s  + \ham_v + \ham_{\mathrm{mix}}.
\end{equation}
The idea is to implement $\ham\ket{\psi}$
by applying these Hamiltonians successively.
For $\ham_s\ket{s}$ we can use the existing \textsc{Fermions} code
with minimal changes.
This is implemented under the function \mintil{_slaterop!()}.

The next term, $\ham_v$ can be written as
\begin{equation}
    \ham_v =
    \sum_{k,\sigma} \epsilon_k \cdag\ks \cc\ks
    - \sum_{\langle ij \rangle, \sigma} \left(
    t_{ij} \cdag_{i\sigma} \cc_{j\sigma}
    + t_{ij}^* \cdag_{j\sigma} \cc_{i\sigma} \right),
    \label{eq:H_vector}
\end{equation}
with energy levels $\epsilon_k$ and hopping amplitudes $t_{ij}$.
$\langle\rangle$ denotes nearest-neighbor hopping.
The expansion by excitation is given by
\begin{equation}
    \ham_v
    \approx \ham_{v0} + \ham_{v1} + \ham_{v2} + \cdots
    = \sum_{i=0}^e \ham_{vi}.
\end{equation}
In the following we will explain each term of the expansion
and how to store and apply it to a hybrid wave function efficiently.
As the Hamiltonian in~\cref{eq:H_vector} conserves the particle number
\begin{equation}
    [\ham_v, \hat{N}] = \sum_{k, \sigma} [\ham_v, \cdag\ks \cc\ks] = 0,
\end{equation}
we can treat each excitation separately.
For the struct it is enough to store the energy levels
and hopping amplitudes inside vectors
($\symbf{\epsilon}$, $\mathbf{t}$ respectively),
since mapping between indices and sites is established
in~\cref{subsec:vector-indexing}.
In the current implementation
we could either supply vectors with
$\dim(\symbf{\epsilon}) = n_v$ and $\dim(\mathbf{t}) = n_v - 1$
internally copying them for both spin orientations,
or $\dim(\symbf{\epsilon}) = 2n_v$ and
$\dim(\mathbf{t}) = 2(n_v - 1)$ which sets spin-dependent energy and hopping
($\epsilon_k \rightarrow \epsilon\ks$ and $t_{ij} \rightarrow t_{ij\sigma}$).
At the moment both vectors are restricted to real values.

Based on the information of $\symbf{\epsilon}$ and $\mathbf{t}$
it is possible to
get the result of any excitation Hamiltonian $\ham_{ve}$ applied
on the vector $\ket{v^e}$ without creating the matrix explicitly
(matrix-free method).
The disadvantage would be that each matrix-vector result must be calculated
on-the-fly,
e.g., for a hybrid representation limited to double excitation (HR-2)
with 100 determinants,
a total of $3\cdot100$ matrix-vector calculations must be done with
the factor 3 resulting from the maximum excitation (including no excitation).
As we want this to be fast, this is not ideal.
We could therefore create a matrix for each excitation once
and cache it if it is feasible memory-wise.
They can then be applied onto each vector
using built in matrix-vector multiplication
which has been optimized for decades already.
In the aforementioned example three matrices would be created.

For no excitation, the hopping terms of $\ham_v$ in \cref{eq:H_vector} can
be ignored since all bath sites are either full ($v_k$) or empty ($c_k$).
The Hamiltonian is therefore just a scalar
\begin{equation}
    \ham_{v0} = \sum_{v_k,\sigma} \epsilon_{v_k} \eqcolon \epsilon_{v0}.
\end{equation}

Given an excitation $e$,
we should first estimate the number of nonzero entries $\nnz$
inside $\ham_{ve}$ to determine whether we could even store the matrix.
If an excitation is at location $(k, \sigma)$,
we must store one value on the diagonal
for the energy $\epsilon_0 \mp \epsilon\ks$
with the sign depending on whether
a valence or conduction bath site is excited,
and two values for nearest-neighbor hopping
$(k, \sigma) \rightarrow (k\pm1, \sigma)$.
This results altogether in
\begin{equation}
    \nnz <
    (1+2e) \binom{2n_{v}}{e}
\end{equation}
with the inequality resulting from the fact that
excitations in the first and last spin-site can only hop in one direction.
Using this result, we can estimate the fill factor of a matrix,
which is defined as the ratio of nonzeros to matrix dimension
\begin{equation}
    f < \frac{\displaystyle(1+2e)\binom{2n_v}{e}}
    {\displaystyle\binom{2n_v}{e}^2}
    \ll 1.
\end{equation}
Since we assume that $e \ll n_v$,
this value is much smaller than one,
which suggests a sparse matrix storage method.
The Julia documentation\footnote{
    \url{https://docs.julialang.org/en/v1/stdlib/SparseArrays/}}
lists the compressed sparse column (CSC) format.
We can set an upper bound of the storage size using the coordinate format (COO)
keeping three \qty{64}{\bit} numbers for each nonzero entry.
As displayed in~\cref{tab:ham-storage},
this poses no problem for singles and doubles,
while triples could be problematic for consumer grade hardware.

\begin{table}[ht]
    \caption{Storage estimation (upper bound) in byte
        of the Hamiltonian in COO format.}
    \label{tab:ham-storage}
    \centering
    \begin{tabular}{l S S S}
        \doubletoprule
                   & \multicolumn{3}{c}{$n_b$ (bath sites)}                 \\
        \cmidrule{2-4}
        Excitation & {100}                                  & {200} & {300} \\
        \midrule
        Single     & 14e3                                   & 29e3  & 43e3  \\
        Double     & 2e6                                    & 10e6  & 22e6  \\
        Triple     & 221e6                                  & 2e9   & 6e9   \\
        \doublebottomrule
    \end{tabular}
\end{table}

For single excitation we can use a more optimized storage format.
As the matrix is tridiagonal, we can use the existing structs
\mintil{Tridiagonal} for complex values
or \mintil{SymTridiagonal} for real values.

Currently, the modified \textsc{Fermions} package features
up to double excitation.
Matrices are cached when creating a new instance of \mintil{CIOperator}.
The matrix-vector multiplication between $\ham_v$ and $\ket{v}$
is implemented inside the function \mintil{_vectorop!()}
using five-argument \mintil{mul!()}.

Lastly, let us turn to $\ham_\mathrm{mix}$.
We can understand the mixed operator as
removing a particle from a spin-site $n_{i\sigma}$ inside $\ket{s}$,
and adding it to spin-site $n_{j\sigma}$ inside $\ket{v}$ or vice versa.
This can be stored inside the struct \mintil{CIOperatorMixed}.
Constructing this object works from tuples $(i, j, t)$
which automatically creates four operators
\begin{equation}
    t \cdag_{i\up} \cc_{j\up}, \quad
    t \cdag_{j\up} \cc_{i\up}, \quad
    t \cdag_{i\dn} \cc_{j\dn}, \quad
    t \cdag_{i\dn} \cc_{j\dn}.
\end{equation}
We can calculate $\ham_{\mathrm{mix}} \ket{\psi}$ by first applying
the corresponding operator on $\ket{s}$ using already
existing infrastructure.
As the sign of the phase depends on the occupation of
all spin-sites between $i$ and $j$ (see~\cref{eq:hopping-1,eq:hopping-2}),
we can split it into two components: a bit phase and vector phase.
For the former we can store a mask which considers
all sites following $(i, \sigma)$,
e.g., if it is at the fourth location,
the mask contains ones from digit five onwards ($\ldots11110000$).
We can interpret adding/removing a particle in $\ket{v}$ as
lowering or raising the excitation and store it as a boolean.
If the excitation inside $\ket{v}$ is lowered,
the single excitation state on $(j, \sigma)$ is mapped to
the state without excitation.
The excitation $c^{(\dagger)}_{j\sigma}$ must be at the first
valence (conduction) bath site.
If the excitation $j$ is at the first valence bath site,
the vector phase is one;
otherwise the excitation is at the first conduction bath site,
resulting in the sign $(-1)^{n_f}$.
For double excitation at $((j, \sigma), (k, \sigma'))$,
the state is mapped to the single excitation state $(k, \sigma')$.
If $(k, \sigma')$ has a smaller index than $(j, \sigma)$,
the aforementioned phase is multiplied by $(-1)$,
since one particle more/less exists.
Raising the excitation can be done similarly.
